\section{Aula 11}

Vimos outra prova do resultado da aula passada,
dessa vez usando uma ferramenta super poderosa. A desigualdade de OSSS.
Na verdade, vimos a versão mais poderosa do Teorema.
\begin{theorem}
	\label{trm:sharp_crossing_transition}
	Seja $R_n = [3n]\times[n]$ e $p \in [0,1]$ tal que $\eta <
	\Pr(H(R_n)) < 1 - \eta$
	para todo $n$. Então, existe $\eps > 0$ tal que
	\[
		\frac{d\Pr_p(H(R_n))}{dp} \geq n^{\eps}.
	\]
\end{theorem}
Lembremos que isso prova automáticamente que $p^{-} = p^+$ e que
$1/2$ é o ponto crítico.

Vamos precisar de algumas definições.
Seja $f : \{0,1\}^n \to \{0,1\}$ uma função booleana do hipercubo $Q_n$ e
como sempre, adicionamos a este a medida de probabilidade em que cada coordenada
é $1$ com probabilidade $p$. Estamos interessados em  $\Pr_p
\{f(\omega) = 1\} = \Ex_p f$.

Agora suponha que temos um algoritmo $A$ que revela as coordenadas de $\omega$
uma por uma e que pode usar as informações anteriores. Mais precisamente,
temos uma sequência de variáveis aleatórias $W_1, \dots, W_n \in [n]$
que representam a $i$-ésima coordenada a ter sido revelada pelo algoritmo.
Pedimos que $W_1 \in \cF_1$ e que $W_{i+1} \in \sigma((W_j,
\omega_j)_{j \leq i})$, isso é,
a sequência, a menos da primeira variável é previsível. Note que é possível que
revelar coordanadas não nos dê nenhuma informação. Por exemplo,
pode ser que o valor de $f$ esteja determinado olhando só para três coordenadas,
$W_1, W_2$ e $W_3$ ($f$ no contexo de percolação poderia ser tem um
caminho de tamanho $3$)
então $W_4 \dots$ não precisariam mais ser observadas, então tomariam
valor arbitrário.

Para ser super exato, definimos um stopping time
\[
	\tau = \inf \{t : \omega'|_{(W_i)\leq t} = \omega|_{(W_i) \leq t}
	\implies f(\omega') = f(\omega)\}
\]
e setamos $W_j = *$ (qualquer coisa que faz sentido) para todo $j > \tau$.
Por fim, dadas $f$, algoritmo $A$ e coordenada $i$ definimos
\[
	\text{Inf}_i(f) = \Pr_p(i \text { pivotal para } f) \quad \& \quad
	\delta_i^A = \Pr_p(\exists j \leq \tau : W_j = i).
\]
Para por em palavras, $\delta_i^A$ é a probabilidade de que durante o
nosso algoritmo,
exploramos a coordenada $i$.

\begin{theorem}[OSSS]
	\label{trm:osss}
	Seja $f$ uma função booleana $f : \{0,1\}^n \to \{0,1\}$ e $A$ um algoritmo
	para determinar $f$. Então,
	\[
		\text{Var}(f) \leq p(1-p) \sum_{i = 1}^n \text{Inf}_i(f) \delta_i^A(f).
	\]
\end{theorem}

\begin{remark}
	Note que um lowerbound em $\text{Var}(f)$ e um upperbound em $\max_i
	\delta_i^A(f)$
	implicam um lowerbound em $\frac{d\Ex_p f}{dp}$ por Margulis Russo. É assim
	que aplicamos a desigualdade de OSSS.
\end{remark}

Vou dar só um sketch da prova, é bem bonitinha.
\begin{proof}
	Notamos primeiramente que $\text{Var}(f) = \Ex (f - \Ex f)^2 = \Ex
	|f(\omega) - f(\omega')|/2$
	onde $\omega$ e $\omega'$ são independentes. A ideia da prova será
	fazer Lindenberg swapping
	com um  coupling esperto, revelando $\omega$ com o algoritmo $A$ e supor que
	as outras coordenadas pertencem a $\tilde{\omega}$. Para isso,
	vamos supor que reordenamos as coordenadas de acordo com $W_1, \dots W_n$ e
	consideramos
	\[f^{(k)} =  f(\omega'_{W_1}, \dots, \omega'_{W_k}, \omega_{W_{k+1}}, \dots,
	\omega_{W_{\tau}}, \omega'_{W_{\tau + 1}}, \dots, \omega'_{W_n})\]

	Segue pela definição de $\tau$ que
	\[
		\Ex|f(\omega) - f(\omega')| \leq \Ex \sum_{k = 0}^{n-1} |f^{(k)} -
		f^{(k+1)}|  = \Ex \sum_{k = 0}^{\tau - 1} |f^{(k)} - f^{(k+1)}|
	\]
	Abrindo a conta, temos
	\[
		\Ex \sum_{k = 0}^{n-1} |f^{(k)} - f^{(k+1)}| = \sum_{k = 0}^{n-1}
		\sum_{i = 1}^n \Ex \id[k - 1 \leq \tau]\id[W_k = i] |f^{(k-1)} - f^{(k)}|.
	\]
	Agora condicionamos em $\cF_{k-1}$ para cada esperança e vemos que temos
	\[
		\Ex \id[k - 1 \leq \tau]\id[W_k = i] |f^{(k-1)} - f^{(k)}| = \Ex
		\id[k-1 \leq \tau] \id[W_{k-1} = i] \Ex[ |f^{(k-1)} - f^{(k)}| | \cF_{k-1}]
	\]
	Agora essa última variável é $1$ se e somente se $\omega'_i \neq
	\omega_i$ e $i$ for pivotal com probabilidades independentes!
	segue que essa esperança é precisamente $2p(1-p)\text{Inf}_i(f)
	\Pr(k - 1 \leq \tau , W_{k-1} = i)$. Somando
	e trocando a ordem temos exatamente
	\[
		2p(1-p)\sum_{k = 0}^{n-1} \sum_{i = 1}^n \text{Inf}_i(f) \Pr(k-1
		\leq \tau, W_{k-1} = i) =
		2p(1-p) \sum_{i = 1}^n \text{Inf}_i(f) \delta_i^A(f).
	\]
\end{proof}

Vamos aplicar isso para provar \ref{trm:sharp_crossing_transition}.
\begin{proof}
	Seja $f = \id[H(R)]$ e $p \in [p^-, p^+]$, logo $\Ex_p f \in (\eta, 1-\eta)$,
	$\text{Var}(f) \geq \eta^2$ e $p(1-p) \geq c$. Temos portanto,
	\[
		\frac{d\Ex_p f}{dp} \geq \frac{c\eta^2}{\max_i \delta_i^A}.
	\]
	Considere o seguinte algoritmo. Escolhemos $X \in [n,2n]$ aleatóriamente
	e exploramos os caminhos saindo da linha vertical ${X} \times [n]$.
	A probabilidade
	de uma aresta $e$ ser revelada é cotada por $\Pr(d(L_X, e) \leq
	\sqrt{n}) + \Pr(e \leftrightarrow \partial B_{\sqrt{n}})$.
	Pela forma que escolhemos $X$, a primeira probabilidade é cotada por
	$O(n^{-1/2})$,
	a segunda é cotada por $O(n^{-\eps})$ por Russo-Seymour-Welsh.
	Isso completa a prova.
\end{proof}
